% ============================================================
%   Chapter 2  --  Project Planning and Management
% ============================================================

\chapter{Project Planning and Management}\label{ch:planning}

The previous chapter framed CALA-SafeRL as a research problem. This chapter frames it as a
\emph{managed engineering project}. It maps the work onto the five canonical project-management
stages (Section~\ref{sec:plan-approach}), decomposes it into a work breakdown structure with an
effort estimate (Section~\ref{sec:plan-wbs}), records the milestones
(Section~\ref{sec:plan-milestones}), the task precedences and the resulting schedule and Gantt
chart (Section~\ref{sec:plan-schedule}), the human-resources plan
(Section~\ref{sec:plan-hr}) and, finally, the risk register and the deviations that actually
occurred (Section~\ref{sec:plan-risk}). The plan presented here is the \emph{baseline} against
which execution was monitored; because the engineering methodology is iterative
(Section~\ref{sec:intro-methodology}), the baseline was re-planned at the end of each
design\,$\to$\,implementation\,$\to$\,training\,$\to$\,analysis loop rather than followed rigidly.

\section{Project-management approach}\label{sec:plan-approach}

The project is organised around the five standard project-management stages (definition,
planning, execution, monitoring and control, and closure), adapted to a research-and-development
context in which the requirements are partly discovered during execution.
Table~\ref{tab:stages} maps each stage to the concrete activities and artefacts of CALA-SafeRL.
A point worth highlighting is the \emph{monitoring and control} stage: in a conventional software
project it is realised through status reports, whereas here it is realised through engineering
instrumentation: the live metrics dashboard (Section~\ref{sec:design-metrics}), the
shield-off probe (Section~\ref{sec:dep-probe}) and the KL and curriculum gates, which provide
continuous, quantitative feedback on whether the work is converging.

\begin{table}[H]
  \centering
  \caption{The five project-management stages mapped to concrete activities and artefacts of
           CALA-SafeRL.}
  \label{tab:stages}
  \small
  \begin{tabular}{@{}lp{6.0cm}p{4.6cm}@{}}
    \toprule
    \textbf{Stage} & \textbf{Activities in this project} & \textbf{Artefacts} \\
    \midrule
    Definition            & Problem framing, research question, scope and objectives
                          & Sections~\ref{sec:intro-context}--\ref{sec:intro-scope} \\
    Planning              & Work breakdown, schedule, resource and risk plans
                          & This chapter \\
    Execution             & Iterative design, implementation, training and analysis loops
                          & Chapters~\ref{ch:design}--\ref{ch:dependence} \\
    Monitoring \& Control & Live metrics dashboard, shield-off probe, KL and curriculum gates,
                            rolling re-planning
                          & Sections~\ref{sec:design-metrics}, \ref{sec:dep-probe} \\
    Closure               & Final evaluation, memoir, defense, reproducibility package
                          & Chapter~\ref{ch:conclusions}, Appendix~\ref{app:reproducibility} \\
    \bottomrule
  \end{tabular}
\end{table}

\section{Work breakdown structure and effort estimation}\label{sec:plan-wbs}

The project is decomposed into eleven work packages (WP0--WP10). WP0 (management) and WP10
(writing) run continuously alongside the technical packages; the remaining packages map directly
onto the specific objectives SO1--SO8 of Section~\ref{sec:intro-objectives}.
Table~\ref{tab:wbs} lists each package, its main tasks, the objective(s) it serves and the
estimated effort. The total of \(550\,\mathrm{h}\) is the author's post-hoc reconstruction of the
effort actually invested; it is consistent with, and somewhat above, the nominal workload of a
\(12\)-ECTS Bachelor's thesis (\(\approx 300\text{--}360\,\mathrm{h}\)), the difference reflecting
the research-heavy nature of the shield-dependence study (WP9) and the large number of training
runs required (WP8).

\begin{table}[H]
  \centering
  \caption{Work breakdown structure and effort estimation. The linked specific objectives (SO)
           are those of Section~\ref{sec:intro-objectives}.}
  \label{tab:wbs}
  \small
  \begin{tabular}{@{}llp{6.0cm}lr@{}}
    \toprule
    \textbf{ID} & \textbf{Work package} & \textbf{Main tasks} & \textbf{Obj.} & \textbf{Hours} \\
    \midrule
    WP0  & Project management        & Planning, scheduling, supervision meetings, risk tracking, version control & -- & 35 \\
    WP1  & Research \& state of the art & RL and Safe-RL literature, shielding, CARLA, bicycle model & -- & 50 \\
    WP2  & Environment \& wrappers   & \texttt{CarlaEnv}, synchronous mode, semantic LIDAR, Waypoint features, 739-d observation & SO1 & 60 \\
    WP3  & PPO agent                 & Actor--Critic network, GAE, clipped surrogate, KL early stop, normalisation & SO2 & 55 \\
    WP4  & Safety shields            & Basic shield, adaptive-horizon shield, bicycle model, shield credit assignment & SO3, SO4 & 70 \\
    WP5  & Reward shaper             & Dense reward components, speed/lane gates, unit tests & SO5 & 30 \\
    WP6  & Curriculum manager        & Five-stage NPC ramp, rollback rule, in-place update & SO6 & 22 \\
    WP7  & Metrics \& dashboard      & SQLite WAL logger, Streamlit/Plotly dashboard & SO7 & 25 \\
    WP8  & Experimentation           & Training runs, shield ablations, evaluation protocol, analysis & -- & 55 \\
    WP9  & Shield-dependence study   & Shield-off probe, LLM meta-controller, behavioural-cloning teacher, failure analysis & SO8 & 68 \\
    WP10 & Thesis writing \& defense & Memoir, figures, reproducibility guide, defense preparation & -- & 80 \\
    \midrule
    \multicolumn{4}{r}{\textbf{Total}} & \textbf{550} \\
    \bottomrule
  \end{tabular}
\end{table}

\section{Milestones}\label{sec:plan-milestones}

Six milestones mark the closure of a major block of work and gate the start of the next.
Table~\ref{tab:milestones} lists them with their target dates and deliverables. Milestones M3--M5
correspond to the three empirical achievements of the project: a working shielded agent, the main
shield ablation, and the resolution of the shield-dependence problem.

\begin{table}[H]
  \centering
  \caption{Project milestones.}
  \label{tab:milestones}
  \small
  \begin{tabular}{@{}llp{7.4cm}@{}}
    \toprule
    \textbf{ID} & \textbf{Target} & \textbf{Milestone / deliverable} \\
    \midrule
    M1 & Sep 2025 & Project definition and baseline plan approved \\
    M2 & Nov 2025 & State of the art consolidated; core bibliography fixed \\
    M3 & Jan 2026 & Full wrapper chain and PPO agent operational; first shielded training run \\
    M4 & Mar 2026 & Main experimentation complete (shield ablation and curriculum) \\
    M5 & May 2026 & Shield-dependence study complete (LLM negative result, BC positive result) \\
    M6 & Jun 2026 & Thesis submitted and defended \\
    \bottomrule
  \end{tabular}
\end{table}

\section{Task scheduling and Gantt chart}\label{sec:plan-schedule}

The technical work packages form a mostly linear dependency chain: the environment (WP2) must
exist before the agent (WP3), both are needed before the shields (WP4), and all of the
components must be in place before systematic experimentation (WP8) and the dependence study
(WP9). Table~\ref{tab:precedence} records the precedences; the resulting \emph{critical path} is
\(\text{WP1}\to\text{WP2}\to\text{WP3}\to\text{WP4}\to\text{WP8}\to\text{WP9}\to\text{WP10}\).
The reward shaper, curriculum manager and metrics pipeline (WP5--WP7) are off the critical path
and were developed in parallel with the shields.

\begin{table}[H]
  \centering
  \caption{Task precedence. WP0 and WP10 run continuously alongside the technical packages.}
  \label{tab:precedence}
  \small
  \begin{tabular}{@{}lll@{}}
    \toprule
    \textbf{WP} & \textbf{Predecessors} & \textbf{Relationship} \\
    \midrule
    WP0  & --                           & Continuous \\
    WP1  & --                           & Start \\
    WP2  & WP1                          & Finish-to-start \\
    WP3  & WP2                          & Finish-to-start \\
    WP4  & WP2, WP3                     & Finish-to-start \\
    WP5  & WP2                          & Finish-to-start (parallel) \\
    WP6  & WP2                          & Finish-to-start (parallel) \\
    WP7  & WP2                          & Finish-to-start (parallel) \\
    WP8  & WP3, WP4, WP5, WP6, WP7      & Finish-to-start \\
    WP9  & WP8                          & Finish-to-start \\
    WP10 & WP1--WP9                     & Continuous \\
    \bottomrule
  \end{tabular}
\end{table}

Figure~\ref{fig:gantt} shows the baseline schedule as a Gantt chart over the ten months of the
project. The diamonds mark the milestones of Table~\ref{tab:milestones}. The chart makes the
front-loading of research (WP1) and the back-loading of writing (WP10) explicit, and shows the
two-month overlap between the close of experimentation (WP8) and the start of the dependence
study (WP9) that occurred in spring 2026.

\begin{figure}[H]
  \centering
  \includegraphics[width=\textwidth]{imgs/gantt_chart}
  \caption{Baseline Gantt chart of the project (September 2025 to June 2026). Bars are the
           work packages of Table~\ref{tab:wbs}; diamonds are the milestones of
           Table~\ref{tab:milestones}. Management (WP0) and writing (WP10) span the whole
           project; the technical critical path runs WP2, WP3, WP4, WP8 and WP9.}
  \label{fig:gantt}
\end{figure}

\section{Human-resources plan}\label{sec:plan-hr}

Although the project was carried out by a single author, the work spanned several professional
\emph{roles}, each with a distinct skill profile and, in the budget of Chapter~\ref{ch:budget}, a
distinct cost rate. Table~\ref{tab:hr-roles} distributes the \(550\,\mathrm{h}\) of
Table~\ref{tab:wbs} across these roles. The RL/ML Engineer role dominates, as expected for a
reinforcement-learning thesis, but the Software Engineer role is also substantial because the
CARLA wrapper, the metrics pipeline and the dashboard are non-trivial systems-engineering tasks.

\begin{table}[H]
  \centering
  \caption{Roles assumed by the author and their effort. The hourly rates are used in the budget
           (Chapter~\ref{ch:budget}).}
  \label{tab:hr-roles}
  \small
  \begin{tabular}{@{}lp{5.8cm}rr@{}}
    \toprule
    \textbf{Role} & \textbf{Main responsibilities} & \textbf{Hours} & \textbf{Rate (€/h)} \\
    \midrule
    Project Manager   & Planning, scheduling, risk management, supervision liaison      & 45  & 40 \\
    RL/ML Engineer    & RL research, PPO, shield logic, reward design, results analysis & 256 & 35 \\
    Software Engineer & CARLA wrapper, metrics pipeline, dashboard, tooling             & 134 & 30 \\
    QA/Test Engineer  & Unit tests, deterministic evaluation, training-run validation   & 45  & 25 \\
    Technical Writer  & Memoir, figures, documentation                                 & 70  & 22 \\
    \midrule
    \multicolumn{2}{@{}l}{\textbf{Total effort}} & \textbf{550} & \\
    \bottomrule
  \end{tabular}
\end{table}

The project also relied on a \emph{support team} that is not part of the author's effort but is
costed in Chapter~\ref{ch:budget}: the project Director, who supervised the work and reviewed the
memoir; the Bachelor-thesis Coordinator, responsible for the administrative process; and the
three-member defense tribunal.

\section{Risk management and deviations}\label{sec:plan-risk}

Table~\ref{tab:risks} is the project risk register. Because the plan is presented retrospectively,
the final column records whether each risk actually materialised. Four of the six risks did, and
the way they were handled is documented in the lessons learned of
Section~\ref{sec:concl-lessons}: the curriculum-change deadlock was resolved by updating the NPC
count in place rather than rebuilding the environment; the \texttt{log\_std} saturation was
resolved by the tanh-squashed parameterisation; the shortage of compute for multi-seed runs was
\emph{accepted} and reported honestly (a single seed with seed-aligned cross-configuration
evaluation); and the late-emerging LLM reward-exposure line was de-scoped to future work to
protect the core objectives. The \(5\%\) contingency reserve in the budget
(Chapter~\ref{ch:budget}) is the financial counterpart of this risk exposure.

\begin{table}[H]
  \centering
  \caption{Risk register. The final column records whether each risk materialised during
           execution.}
  \label{tab:risks}
  \small
  \begin{tabular}{@{}p{4.4cm}ccp{4.4cm}c@{}}
    \toprule
    \textbf{Risk} & \textbf{Lik.} & \textbf{Imp.} & \textbf{Mitigation} & \textbf{Status} \\
    \midrule
    CARLA server instability / crashes               & High & Med  & Synchronous headless mode; checkpoint every 200 episodes & Managed \\
    Deadlock when rebuilding the env on curriculum change & Med  & High & In-place \texttt{num\_npc} update, no reconnection & Resolved \\
    Training instability (\texttt{log\_std} saturation)   & Med  & High & tanh-squashed \texttt{log\_std} parameterisation & Resolved \\
    Insufficient compute for multi-seed runs         & High & Med  & Seed-aligned cross-configuration evaluation; honest reporting & Accepted \\
    Scope creep (LLM reward-exposure line)           & Med  & Med  & De-scoped to future work; core objectives protected & De-scoped \\
    Single-workstation hardware failure              & Low  & High & Frequent checkpointing; Git version control & Not hit \\
    \bottomrule
  \end{tabular}
\end{table}

With the schedule, resources and risks established, Chapter~\ref{ch:budget} turns this plan into a
cost estimate.

\clearpage
